\documentclass[10pt]{article}

\usepackage[utf8]{inputenc}

\usepackage[T1]{fontenc}

\usepackage{amsmath}

\usepackage{amsfonts}

\usepackage{amssymb}

\usepackage[version=4]{mhchem}

\usepackage{stmaryrd}



\begin{document}

$\gamma_{v},-\infty<\gamma_{v} \leqslant+\infty$, наз. Робена постолнной р. п. $\bar{R}_{v}$, $c_{v}=e^{-\nu_{v}}$ есть емкость края $\partial \bar{R}_{v}$ (относительно фиксированного полюса $\left.p_{0} \in R\right)$. При стремлении $v$ к $\infty$ значения $G_{v}\left(p, p_{0}\right)$ и $p_{v}$ могут только возрастать. Функция Грина открытой р. п. $\boldsymbol{R}$ определяется как предел $G\left(p, p_{0}\right)$ возрастающей последовательности $\left\{G_{v}\left(p, p_{0}\right)\right\}$, если он существует; в противном случае, когда



\[

\lim _{v \rightarrow \infty} G_{v}\left(p, p_{0}\right) \equiv+\infty,

\]



говорят, что р. п. $R$ не имеет функции Грина, причем существование или несуществование функции Грина не зависит от выбора полюса $p_{0} \in R$. Класс р. п., для к-рых функция Грина не существует, обозначается $\widehat{G}_{G}$. Иными словами, класс $6_{G}$ характеризуется тем, что



\[

\lim _{v \rightarrow \infty} \gamma_{v}=+\infty \text { или } \lim _{v \rightarrow \infty} c_{v}=0

\]



шричем эти соотношения также не зависят от выбора полюса.



Пусть $R$ - открытая р. П. и $\left\{\Delta_{v}\right\}_{v=1}^{\infty}$ - т. н. о п р едел я ющая послед о вательность замкнутых на $R$ областей $\Delta_{v}$, т. е. такая последовательность, что 1) граница $\Delta_{v}$ есть простая замкнутая кривая на $R$; 2) $\Delta_{v+1} \subset \Delta_{v}, v=1,2, \ldots$; 3) $\bigcap_{v=1}^{\infty} \Delta_{v}=\varnothing$, то есть $\Delta_{v}$ не компактны на $R$. Две определяющие последовательности $\left\{\Delta_{v}\right\}$ и $\left\{\Delta_{v}^{\prime}\right\}$ әквивалентны, если каждому $v$ соответствуют такие $n$ и $m$, что $\Delta_{n}^{\prime} \subset \Delta_{v}$ и $\Delta_{m} \subset \Delta_{v}^{\prime}$. Классы эквивалентности определяющих последовательностей наз. г р а ни ч ным и әле ментам и р. п. $R$, a совокупность всех граничных әлементов образует и де а ль н ую г раниц у $\Gamma$ р. п. $R$, рассматриваемой как топологич. поверхность. Напр., идеальная граница единичного круга $D$ состоит из одного граничного элемента. Отметим, что функция Грина открытой p. п. $R$, в отличие от случая гиперболич. конечной р. п., не обязательно обращается в нуль на всех әлементах идеальной границы $Г$. Класс $6_{G}$ характеризуется также как класс р. п. с идеальной границей нулевой емкости, или, короче, как класс р. п. с вулевой границей; если $R \notin G_{G}$, то $\lim _{v \rightarrow \infty} c_{v}=c>0$ наз. емкостью идеальной границы. Существование или несуществование функции Грина р. п. $R$, а также объем других функциональных классов на $R$ определяются прежде всего этой и другими более тонкими характеристиками идеальной границы, связанными с самими функциональными классами.



Основным и функци ональны ми к ла сс а м и $W$ на р. п. $R$ являются следующие:



$A B$ - класс ограниченных однозначных аналитич. функций на $R$;



$A D$ - класс однозначных аналитич. Функций $f(z)$ на $R$ с конечным Дирихле интегралом



\[

D_{R}(f)=\iint_{R}\left|\frac{d w}{d z}\right|^{2} d x d y, z=x+i y ;

\]



$H P, H B$ и $H D-$ классы однозначных гармонич. функций на $R$ соответственно положительных, ограниченных и с конечным интегралом Дирихле. Эти классы могут комбинироваться, напр. $A B D$ - класс ограниченных однозначных аналитич. функций на $R$ с конечным интегралом Дирихле. Для соответствующих классов $\sigma_{W}$ р. п. $R$ установлены следующие строгие включения и равенства:



\[

\sigma_{G} \subset 6_{H P} \subset 6_{H B} \subset \sigma_{A B} \subset \sigma_{A B D}=6_{A D},

\][



\textbackslash sigma\_\{H B\} \textbackslash subset \textbackslash sigma\_\{H D\} \textbackslash subset \textbackslash hat\{\textbackslash sigma\}\textit{\{A D\}, \textbackslash hat\{\textbackslash sigma\}}\{H B D\}=\textbackslash sigma\_\{H D\}

]



Для плоских областей $R$ эти соотношения упрощаются:



\[

\begin{gathered}

\sigma_{G}=\sigma_{H P}=\sigma_{H B}=\sigma_{H D}, \\

\sigma_{A B} \subset \sigma_{A B D}=\sigma_{A D} .

\end{gathered}

\]



Важное значение имеют также к л а с с ы $\mathrm{X}$ а $\mathrm{p}$ д и $A H_{p}, 0<p \leqslant+\infty$, однозначных аналитич. функций $w=$ $=f(z)$ на р. п. $R$. При $0<p<+\infty$ функция $f \in A H_{p}$, если субгармонич. функция $|f|^{p}$ имеет на всөй р. п. $R$ гармонич. мажоранту, а $A H_{\infty}=A B$ (см. Граничные свойства аналитических функций).



Р. п. параболич. типа принадлежит классу $\mathfrak{\sigma}_{G}$, поэтому иногда вопрос о характеризации р. п. класса $\widehat{\sigma}_{G}$ наз. обобще нн о й пробле мо й ти па. Имеется много результатов, в к-рых в различных терминах устанавливаются условия принадлежности р. п. указанным выше классам. Глубокие исследования были посвящены выяснению внутренних свойств р. п. определенных классов. В частности, оказалось, что р. п. с нулевой границей во многих отношениях аналогичны замкнутым р. п. На них строятся аналоги абелевых дифференциалов и соответствующие интегралы.



Более тонкие свойства идеальной границы р. п. $R$ удается исследовать также с помотью различных компактификаций $R$. Напр., пусть $N(R)$ - винеровская алгебра функций $u$ на р. п. $R$, ограаниченных, непрерывных и гармонизуемых на $R$; последнее означает, что для любой регулярной области $G \subset R$ существует обобщенное решение задачи Дирихле в смысле Винера - Перрона (см. Перрона метод) с граничными данными $u$ на границе $\partial G$. К о м п а к т и фо и ка ц и е ї В и н е р a p. п. $R$ наз. компактное хаусдорфово пространство $R^{*}$ такое, что $R$ есть открытое плотное подпространство $R^{*}$, каждая функция $u \in N(R)$ непрерывно продолжается на $R^{*}$ и $N(R)$ разделяет точки $R^{*}$. Компактификация Винера существует для любой р. п. $R$. Множество $\Gamma(R)=R^{*} \backslash R$ наз. в и н е р о в с к о й и д е а ль ь о й $\mathbf{⿴ 囗}$ а н и це й р. п. $R$, а подмножество $\Delta(R) \subset \Gamma(R)$, состоящее из тех точек $R^{*}$, в к-рых все потенциалы из $N(R)$ обращаются в нуль, - в и н еровск ой гармоничесской г В этих терминах, напр., включение $R \in G_{G}$ равносильно равенству $\Delta(R)=\varnothing$; отсюда получается также строгое включение $\widehat{O}_{H P} \subset \widehat{\vartheta}_{H B}$.



С Р. П. к. связаш также вопрос об устранимых множествах на р. п. Так, компакт $K$ на р. ■. $R$ наз. $A B-$ у с т р а н и м ы м, если для нек-рой окрестности $U \supset$ $\supset K$ на $R$ все $A B$-функции на $\vec{U} \backslash K$ имеют аналитич. продолжение на всю окрестность $U$.



Внимание многих исследователей привлечено также к вопросам классификации римановых многообразий произвольных размерностей $N \geqslant 2$, связанной с рассмотрением описанных выше классов функций.



Лит. см. при ст. Риманова поверхность.



РИМАНОВЫХ ПОВЕРХНОСТЕЙ КОНФОРМНЫЕ КЛАССЫ - классы, состоящие из конформно эквивалентных римановых поверхностей. Замкнутые римановы поверхности (р. п.) имеют простой топологич. инвариант - род $g$; при этом любые две поверхности одного рода гомеоморфны. В простейших случаях топологич. эквивалентность двух р. п. обеспечивает и их принадлежность к одному и тому же Р. п.к. к., то есть конформную эквивалентность, или, говоря иначе, совпадение их конформных структур. Это выполняется, напр., для поверхностей рода 0 , то есть гомеоморфных сфере. В общем случае дело обстоит не так. Еще Б. Риман (B. Riemann) заметил, что классы конформной эквивалентности р. п. рода $g>1$ зависят от $3 g-3$ комплексных параметров, называемых модуляли римановых поверхностей; для конформно эквивалентных поверхностей эти модули совпадают. Случай $g=1$ описан ниже. Если же рассматривать компактные р. п. рода $g$ с $n$





\end{document}